% GENERATED FILE. Edit canonical lessons, then run npm run generate:books.
% canonical-source-hash: fnv1a64:a7c94f66f66034b7
% canonical-lessons: KA-C351-karyacarane, KA-C351-svayamsevaka, KA-C351-appalisu, KA-C351-avarisu, KA-C351-koccihogu

\chapter{The Rescue Operation}
\label{ch:ka-the-rescue-operation}
\hlchaptermodality{351}

\begin{chapteropening}
\textbf{By the end of this chapter:} \emph{I can say a mission, a volunteer, to crash against, to engulf and to be washed away.}

Complete the last lesson of chapter 351: I can say a mission, a volunteer, to crash against, to engulf and to be washed away.
\end{chapteropening}

\section[\texorpdfstring{\kn{ಕಾರ್ಯಾಚರಣೆ} (kāryācaraṇe)}{ಕಾರ್ಯಾಚರಣೆ (kāryācaraṇe)}]{\kn{ಕಾರ್ಯಾಚರಣೆ} (kāryācaraṇe) — a mission; an operation}

\label{lesson:KA-C351-karyacarane}



\begin{quote}
Before the new one: say the Kannada for fame, then the Kannada for to inaugurate.
\end{quote}

\subsection*{You'll want to know: \kn{ಕಾರ್ಯಾಚರಣೆ}}
\textbf{\kn{ಕಾರ್ಯಾಚರಣೆ}} — \emph{kāryācaraṇe} — \textquotedblleft{}a mission; an operation\textquotedblright{}.

\begin{quote}
\kn{ಪ್ರವಾಹದಲ್ಲಿ} \kn{ರಕ್ಷಣಾ} \kn{ಕಾರ್ಯಾಚರಣೆ} \kn{ಎರಡು} \kn{ದಿನ} \kn{ನಡೆಯಿತು}.
\end{quote}

\begin{quote}
\emph{pravāhadalli rakṣaṇā kāryācaraṇe eraḍu dina naḍeyitu.}
\end{quote}

\begin{quote}
The rescue operation in the flood went on for two days.
\end{quote}

Used of rescue, police and army work: \kn{ರಕ್ಷಣಾ} \kn{ಕಾರ್ಯಾಚರಣೆ}, a rescue operation. In a flood it brings food and \kn{ಔಷಧ}, medicine, to people cut off by the water.

\begin{cousinweb}
\textbf{\kn{ಕಾರ್ಯಾಚರಣೆ}} (\emph{kāryācaraṇe}) is borrowed from Sanskrit, a modern compound of \emph{kārya}, \textquotedblleft{}a thing to be done, business\textquotedblright{} (CDIAL 3078), and \emph{ācaraṇa}, \textquotedblleft{}undertaking, performing\textquotedblright{} (Monier-Williams).
\end{cousinweb}

\begin{grammarlens}[title={What you've built}]
One more word for telling what happened.
\end{grammarlens}

\subsection*{Your turn}
\begin{itemize}
\raggedright
  \item \emph{Say it:} \emph{kāryācaraṇe}
  \item \emph{Say it:} \emph{kāryācaraṇe}, once more
  \item \emph{Say it:} \emph{pravāhadalli rakṣaṇā kāryācaraṇe eraḍu dina naḍeyitu.} — The rescue operation in the flood went on for two days.
  \item \emph{Recall:} say the Kannada for fame, then the Kannada for to inaugurate, then say \emph{kāryācaraṇe} again
\end{itemize}

\begin{tcolorbox}[breakable,colback=teal!4,colframe=teal!35!black,title={Before you move on}]
What does \kn{ಕಾರ್ಯಾಚರಣೆ} mean? (\textquotedblleft{}A mission; an operation\textquotedblright{}.) Say it once more.
\end{tcolorbox}
\section[\texorpdfstring{\kn{ಸ್ವಯಂಸೇವಕ} (svayaṃsēvaka)}{ಸ್ವಯಂಸೇವಕ (svayaṃsēvaka)}]{\kn{ಸ್ವಯಂಸೇವಕ} (svayaṃsēvaka) — a volunteer}

\label{lesson:KA-C351-svayamsevaka}



\begin{quote}
Before the new one: say the Kannada for to inaugurate, then the Kannada for a mission.
\end{quote}

\subsection*{You'll want to know: \kn{ಸ್ವಯಂಸೇವಕ}}
\textbf{\kn{ಸ್ವಯಂಸೇವಕ}} — \emph{svayaṃsēvaka} — \textquotedblleft{}a volunteer\textquotedblright{}.

\begin{quote}
\kn{ನೂರಾರು} \kn{ಸ್ವಯಂಸೇವಕರು} \kn{ಸಂತ್ರಸ್ತರಿಗೆ} \kn{ಊಟ} \kn{ತಂದರು}.
\end{quote}

\begin{quote}
\emph{nūrāru svayaṃsēvakaru santrastarige ūṭa tandaru.}
\end{quote}

\begin{quote}
Hundreds of volunteers brought food to the victims.
\end{quote}

Plural: \kn{ಸ್ವಯಂಸೇವಕರು}, \emph{svayaṃsēvakaru}.

\kn{ಸ್ವಯಂ}, oneself, and \kn{ಸೇವಕ}, one who serves: one who serves of his own will.

\begin{cousinweb}
\textbf{\kn{ಸ್ವಯಂಸೇವಕ}} (\emph{svayaṃsēvaka}) is borrowed from Sanskrit, a modern compound of \emph{svayam}, \textquotedblleft{}self, spontaneously\textquotedblright{} (CDIAL 13906), and \emph{sevaka}, \textquotedblleft{}one who serves\textquotedblright{} (Monier-Williams).
\end{cousinweb}

\begin{grammarlens}[title={What you've built}]
One more word for telling what happened.
\end{grammarlens}

\subsection*{Your turn}
\begin{itemize}
\raggedright
  \item \emph{Say it:} \emph{svayaṃsēvaka}
  \item \emph{Say it:} \emph{svayaṃsēvaka}, once more
  \item \emph{Say it:} \emph{nūrāru svayaṃsēvakaru santrastarige ūṭa tandaru.} — Hundreds of volunteers brought food to the victims.
  \item \emph{Recall:} say the Kannada for to inaugurate, then the Kannada for a mission, then say \emph{svayaṃsēvaka} again
\end{itemize}

\begin{tcolorbox}[breakable,colback=teal!4,colframe=teal!35!black,title={Before you move on}]
What does \kn{ಸ್ವಯಂಸೇವಕ} mean? (\textquotedblleft{}A volunteer\textquotedblright{}.) Say it once more.
\end{tcolorbox}
\section[\texorpdfstring{\kn{ಅಪ್ಪಳಿಸು} (appaḷisu)}{ಅಪ್ಪಳಿಸು (appaḷisu)}]{\kn{ಅಪ್ಪಳಿಸು} (appaḷisu) — to crash against; to strike}

\label{lesson:KA-C351-appalisu}



\begin{quote}
Before the new one: say the Kannada for a mission, then the Kannada for a volunteer.
\end{quote}

\subsection*{You'll want to know: \kn{ಅಪ್ಪಳಿಸು}}
\textbf{\kn{ಅಪ್ಪಳಿಸು}} — \emph{appaḷisu} — \textquotedblleft{}to crash against; to strike\textquotedblright{}.

\begin{quote}
\kn{ದೊಡ್ಡ} \kn{ಅಲೆಗಳು} \kn{ದಡಕ್ಕೆ} \kn{ಅಪ್ಪಳಿಸಿದವು}.
\end{quote}

\begin{quote}
\emph{doḍḍa alegaḷu daḍakke appaḷisidavu.}
\end{quote}

\begin{quote}
Huge waves crashed against the shore.
\end{quote}

Past (it …): \kn{ಅಪ್ಪಳಿಸಿತು}, \emph{appaḷisitu}.

What is hit takes -\kn{ಗೆ}. Storms, waves and news \kn{ಅಪ್ಪಳಿಸು}.

\begin{cousinweb}
\textbf{\kn{ಅಪ್ಪಳಿಸು}} (\emph{appaḷisu}) is native Dravidian: DEDR 157 lists Kannada \emph{appaḷisu}, \textquotedblleft{}to strike against, to slap\textquotedblright{}.
\end{cousinweb}

\begin{grammarlens}[title={What you've built}]
One more word for telling what happened.
\end{grammarlens}

\subsection*{Your turn}
\begin{itemize}
\raggedright
  \item \emph{Say it:} \emph{appaḷisu}
  \item \emph{Say it:} \emph{appaḷisu}, once more
  \item \emph{Say it:} \emph{doḍḍa alegaḷu daḍakke appaḷisidavu.} — Huge waves crashed against the shore.
  \item \emph{Recall:} say the Kannada for a mission, then the Kannada for a volunteer, then say \emph{appaḷisu} again
\end{itemize}

\begin{tcolorbox}[breakable,colback=teal!4,colframe=teal!35!black,title={Before you move on}]
What does \kn{ಅಪ್ಪಳಿಸು} mean? (\textquotedblleft{}To crash against; to strike\textquotedblright{}.) Say it once more.
\end{tcolorbox}
\section[\texorpdfstring{\kn{ಆವರಿಸು} (āvarisu)}{ಆವರಿಸು (āvarisu)}]{\kn{ಆವರಿಸು} (āvarisu) — to engulf; to cover}

\label{lesson:KA-C351-avarisu}



\begin{quote}
Before the new one: say the Kannada for a volunteer, then the Kannada for to crash against.
\end{quote}

\subsection*{You'll want to know: \kn{ಆವರಿಸು}}
\textbf{\kn{ಆವರಿಸು}} — \emph{āvarisu} — \textquotedblleft{}to engulf; to cover\textquotedblright{}.

\begin{quote}
\kn{ಬೆಂಕಿಯ} \kn{ಹೊಗೆ} \kn{ಇಡೀ} \kn{ಬೀದಿಯನ್ನು} \kn{ಆವರಿಸಿತು}.
\end{quote}

\begin{quote}
\emph{beṅkiya hoge iḍī bīdiyannu āvarisitu.}
\end{quote}

\begin{quote}
Smoke from the fire covered the whole street.
\end{quote}

Past (it …): \kn{ಆವರಿಸಿತು}, \emph{āvarisitu}.

Smoke, darkness, fear and silence \kn{ಆವರಿಸು}: \kn{ಕತ್ತಲೆ} \kn{ಆವರಿಸಿತು}, darkness fell.

\begin{cousinweb}
\textbf{\kn{ಆವರಿಸು}} (\emph{āvarisu}) is built on Sanskrit \emph{āvaraṇa}, \textquotedblleft{}covering, obstruction\textquotedblright{} (CDIAL 1413), with the Kannada verb ending -\emph{isu}.
\end{cousinweb}

\begin{grammarlens}[title={What you've built}]
One more word for telling what happened.
\end{grammarlens}

\subsection*{Your turn}
\begin{itemize}
\raggedright
  \item \emph{Say it:} \emph{āvarisu}
  \item \emph{Say it:} \emph{āvarisu}, once more
  \item \emph{Say it:} \emph{beṅkiya hoge iḍī bīdiyannu āvarisitu.} — Smoke from the fire covered the whole street.
  \item \emph{Recall:} say the Kannada for a volunteer, then the Kannada for to crash against, then say \emph{āvarisu} again
\end{itemize}

\begin{tcolorbox}[breakable,colback=teal!4,colframe=teal!35!black,title={Before you move on}]
What does \kn{ಆವರಿಸು} mean? (\textquotedblleft{}To engulf; to cover\textquotedblright{}.) Say it once more.
\end{tcolorbox}
\section[\texorpdfstring{\kn{ಕೊಚ್ಚಿಹೋಗು} (koccihōgu)}{ಕೊಚ್ಚಿಹೋಗು (koccihōgu)}]{\kn{ಕೊಚ್ಚಿಹೋಗು} (koccihōgu) — to be washed away}

\label{lesson:KA-C351-koccihogu}



\begin{quote}
Before the new one: say the Kannada for to crash against, then the Kannada for to engulf.
\end{quote}

\subsection*{You'll want to know: \kn{ಕೊಚ್ಚಿಹೋಗು}}
\textbf{\kn{ಕೊಚ್ಚಿಹೋಗು}} — \emph{koccihōgu} — \textquotedblleft{}to be washed away\textquotedblright{}.

\begin{quote}
\kn{ಪ್ರವಾಹದಲ್ಲಿ} \kn{ಹಳೆಯ} \kn{ಸೇತುವೆ} \kn{ಕೊಚ್ಚಿಹೋಯಿತು}.
\end{quote}

\begin{quote}
\emph{pravāhadalli haḷeya sētuve koccihōyitu.}
\end{quote}

\begin{quote}
The old bridge was washed away in the flood.
\end{quote}

Past (it …): \kn{ಕೊಚ್ಚಿಹೋಯಿತು}, \emph{koccihōyitu}.

Things carried off by water: \kn{ಮನೆಗಳು} \kn{ಕೊಚ್ಚಿಹೋದವು}, houses were swept away.

\begin{cousinweb}
\textbf{\kn{ಕೊಚ್ಚಿಹೋಗು}} (\emph{koccihōgu}) is a Kannada compound of \kn{ಕೊಚ್ಚಿ} and \emph{hōgu}, \textquotedblleft{}to go\textquotedblright{} (DEDR 4572); DEDR 1859 lists a Kannada \emph{koccu}, \textquotedblleft{}to cut away\textquotedblright{}, but whether \kn{ಕೊಚ್ಚಿ} is that verb is not settled.
\end{cousinweb}

\begin{grammarlens}[title={What you've built}]
One more word for telling what happened.
\end{grammarlens}

\subsection*{Your turn}
\begin{itemize}
\raggedright
  \item \emph{Say it:} \emph{koccihōgu}
  \item \emph{Say it:} \emph{koccihōgu}, once more
  \item \emph{Say it:} \emph{pravāhadalli haḷeya sētuve koccihōyitu.} — The old bridge was washed away in the flood.
  \item \emph{Recall:} say the Kannada for to crash against, then the Kannada for to engulf, then say \emph{koccihōgu} again
\end{itemize}

\begin{tcolorbox}[breakable,colback=teal!4,colframe=teal!35!black,title={Before you move on}]
What does \kn{ಕೊಚ್ಚಿಹೋಗು} mean? (\textquotedblleft{}To be washed away\textquotedblright{}.) Say it once more.
\end{tcolorbox}
